\documentclass[10pt,twocolumn,letterpaper]{article}

\usepackage{cvpr}
\usepackage{times}
\usepackage{epsfig}
\usepackage{graphicx}
\usepackage{amsmath}
\usepackage{amssymb}
\usepackage{dblfloatfix}
\usepackage{placeins}
\usepackage{fancyhdr}

\usepackage[breaklinks=true,bookmarks=false]{hyperref}

\cvprfinalcopy

\def\cvprPaperID{****}
\def\httilde{\mbox{\tt\raisebox{-.5ex}{\symbol{126}}}}

\ifcvprfinal
\pagestyle{fancy}
\fancyhf{}
\fancyfoot[C]{\thepage}
\renewcommand{\headrulewidth}{0pt}
\renewcommand{\footrulewidth}{0pt}
\fi

\begin{document}

\title{Parallel Logic Expert with Operator Routing}

\author{
Steven Cleasby-Mayeda\\
Oregon State University\\
{\tt\small cleasbys@oregonstate.edu}
{\tt\small steven.cleasby.mayeda@gmail.com}
}
\date{AI535 Deep Learning \\ 25 March 2026}

\maketitle
\begin{abstract}
We introduce a logic-augmented transformer pathway that routes token-level
representations into fixed fuzzy logical operators (AND/OR/NOT) or straight-through-estimation,
so the network learns compositional structure rather than an ensemble of expert networks
with shared architecture design. At each backbone layer, the hidden states are
projected to a lower dimensional K, Q, which the running logic state queries
(cross-attention). The attention matrix is routed, and local
windows over either token or gate-channel axes are composed.
The final logic state is fused with the backbone features before classification.
The proposed architecture augmentation imposes structured, interpretable compositions over
token features. A key design dimension is the axis of composition:
intra-token compositions (feature axis) compares features within a token, while
inter-token compositions (sequence axis) compares features across neighboring
tokens. We evaluate these design choices under controlled settings
with matched parameter budgets, report task performance, mechanism-level
visualizations of the activations of the novel components.
\end{abstract}

\section{Introduction}

\begin{figure*}[tbp]
\centering
\includegraphics[width=0.75\textwidth]{plots/arch1.png}
\caption{Our architecture, backbone hidden-states feed to logic expert, and fused back with
the backbone features before classification}
\label{fig:architecture}
\end{figure*}

Transformers provide strong performance on a wide variety of tasks \cite{vaswani2023attentionneed},
both chain-of-thought (CoT) prompting and reasoning models can greatly
improve performance on reasoning tasks
\cite{wei2023chainofthoughtpromptingelicitsreasoning,Guo_2025}. However,
many CoT use N-shot prompting. Additionally, both CoT and reasoning models
utilize reasoning by iterative, and extended inference, which requires far more
tokens. This increases inference costs for many tasks
\cite{wei2023chainofthoughtpromptingelicitsreasoning}.
Our approach instead introduces a novel architecture for
structured computation within the latent space, no additional
tokens needed. We will test how this constrained, logic-inspired
pathway affects performance on reasoning problems and generalization.

Mixture-of-Experts (MoE) models
\cite{6797059,shazeer2017outrageouslylargeneuralnetworks,fedus2022switchtransformersscalingtrillion}
introduce conditional computation by routing tokens to specialized subnetworks, or weighting
the results of several/all experts.
Our method introduce conditional computation by routing, weighting
tokens to logical operators. In this formulation,
the operator routing is learned rather than expert weights.
When tied with attention, this method may allow robust routing
in the logic space. This space has forced sparsity due to softmax
and local compositions (local intra-token or inter-token windows),
which may allow emergence of specialized operator-feature connections.

Our activation metrics allows learned logic compositions to be interpreted.
Routing entropy and fusion $\alpha$ parameters allow us to measure the
utilization of the logic expert. Experiments on logical
reasoning benchmarks compare the Baseline, the No-Gate Control network
(a baseline with matched parameter count for fair comparison), and the
logic expert variants.
A core design question is how the local operator windows should be applied in the
gate compositions. Our first implementation applies capacity on the channel axis,
(neighboring channels on one token). The next design applies
capacity on the token axis, each channel compares the same indexed feature on
neighboring tokens (channels remain independent). This axis change can have
a large impact on which feature relationships are learnable, which has
large implications for performance.
Figure~\ref{fig:architecture} provides a high-level view of the
architecture.

Contributions of this paper are: (1) inter-token compositional routing with
operator reductions and residual logic-state updates; (2) controlled ablations
over logic dimension, gate count, window size/axis, gate mode, and fusion
$\alpha$; and (3) evaluation on reasoning and transfer settings with
mechanism-level metrics (routing entropy, fusion $\alpha$, and average gradient norm).

\section{Related Work}

Logical and logic-inspired neural architectures span a wide variety of designs,
from classic weightless systems to modern differentiable operator networks
\cite{Aleksander2009ABI}.
Some architectures utilize explicit logical, or logic-inspired
differentiable operators \cite{petersen2022deepdifferentiablelogicgate} as an
alternative to dense transformations, and they provide evidence that structured
logical computation can be more efficient and interpretable than typical
transformers. Our approach builds on this motivation, but integrates operator routing
directly into transformer hidden-state rather than standalone units.

Deep Differentiable Logic Gate Networks combine real-valued logical operators
with continuous relaxations (differentiable approximations of discrete
logical gates used in gradient descent) and can achieve very efficient inference
(one million+ MNIST images/s on a single CPU core)
\cite{petersen2022deepdifferentiablelogicgate}.
Logical Neural Networks (LNNs) \cite{riegel2020logicalneuralnetworks}
map neural units to components of weighted logical formulas, yielding highly
interpretable representations and supporting logic-consistent reasoning.
LNNs use explicitly logic-structured model semantics.
Differentiable Logic Machines (DLMs) \cite{zimmer2023differentiablelogicmachines}
define a continuous relaxation of first-order logic program spaces (via
predicate-weighted parameterization) and learn interpretable logic programs
with gradient descent. In contrast to these approaches, which operate over
structured relational inputs or operations on the input space,
our method integrates logical composition in the feature space,
it learns logical compositions over token features.

MoE architectures use a learned gating mechanism to route
inputs to specialized expert subnetworks. In adaptive mixtures of local experts,
the objective for each expert is decoupled from other experts, enabling
specialization over different input regions, and thus allow both specialized and
general networks \cite{6797059}. These models can also improve inference efficiency
via forced sparse activations.
Conditional computation methods such as MoE and sparse
routing motivates learned token-dependent routing decisions
\cite{shazeer2017outrageouslylargeneuralnetworks,fedus2022switchtransformersscalingtrillion}.
Our model is closest to this line of work, but uses compact differentiable
operators instead of expert-token dispatch among independent expert sub-networks.

Fuzzy-set theory provides continuous, differentiable functions \cite{ZADEH1965338},
We use fuzzy-set operations to estimate logical operations,
which allows training with any off-the-shelf gradient-based optimizer.
Our architecture is contrasted by Neural-Symbolic networks,
where explicit logic is composed in the non-compressed input space
(for example, decision trees or markov-networks)
\cite{besold2017neuralsymboliclearningreasoningsurvey}.

\section{Methods}

\subsection{Architecture and State Representation}

The logic-augmented architecture introduces the logic-expert pathway parallel to
the transformer backbone. At each layer $t$, the backbone produces hidden states:
\begin{equation}
H_t \in \mathbb{R}^{B \times S \times H}.
\end{equation}
In parallel, the logic expert maintains logic states which are updated at every layer:
\begin{equation}
L_t \in \mathbb{R}^{B \times S \times L},
\end{equation}

At each layer, the logic state attends to the backbone via multi-head cross-attention,
producing attended logic features:
\begin{equation}
\tilde{L}_t.
\end{equation}

After the final layer, the logic state is fused with the backbone:
\begin{equation}
H_{\text{fused}} = H_{\text{last}} + \alpha W_f L_{\text{last}},
\label{eq:fusion}
\end{equation}
where $\alpha$ controls the contribution of the logic pathway.

\paragraph{Notation.}
$\pi_t$: routing weights; $s_t$: gate activations; $u_t$: routed gate activations;
$B$: batch size; $S$: sequence length; $H$: backbone hidden dimension;
$L$: logic-state dimension; $G$: gates per operator family (AND/OR/NOT), with
total routed channels $3G$; $C_t$: local window size at layer $t$;
$\tau$: routing softmax temperature.

\subsection{Cross-Attention into Logic Space}

At each layer, logic tokens act as queries, while backbone hidden states are projected
to keys and values. This produces attended logic features:
\begin{equation}
\tilde{L}_t := \text{CrossAttention}(L_{t-1}, H_t).
\end{equation}

These features serve as the input to the routing and composition stage.

\subsection{Routing and Gate Computation}

Routing and gate activation use separate learned projection matrices: $W_r$
produces routing logits, while $W_g$ produces gate activations.

Routings are computed from $\tilde{L}_t$:
\begin{equation}
R_t := \tilde{L}_t W_r, \quad W_r \in \mathbb{R}^{L \times 3G}.
\end{equation}

Routing weights are:
\begin{equation}
\pi_t[b,s,j] = \frac{\exp(R_t[b,s,j]/\tau)}{\sum_{k=1}^{3G} \exp(R_t[b,s,k]/\tau)}.
\end{equation}

Gate activations are computed as:
\begin{equation}
s_t := \sigma(\tilde{L}_t W_g), \quad W_g \in \mathbb{R}^{L \times 3G}.
\end{equation}

The routed soft activations are:
\begin{equation}
u_t^{\mathrm{soft}} := \pi_t \odot s_t.
\end{equation}
In soft mode, $u_t := u_t^{\mathrm{soft}}$. In STE mode (forward binarized), we use
\begin{equation}
u_t^{\mathrm{hard}} := \mathbf{1}[u_t^{\mathrm{soft}} \ge 0.5], \quad
u_t := u_t^{\mathrm{soft}} + \operatorname{sg}(u_t^{\mathrm{hard}} - u_t^{\mathrm{soft}}).
\end{equation}

We split:
\begin{equation}
[u_t^{\text{AND}}, u_t^{\text{OR}}, u_t^{\text{NOT}}] := \operatorname{split}(u_t; G, G, G).
\end{equation}

\subsection{Windowed Operator Composition}

Let $v \in [0,1]^{B \times S \times G}$ denote one operator group.

Windows are causal and defined along either:
inter-token axis (sequence), or intra-token axis (channels).

For inter-token composition:
\begin{equation}
v_{\text{win}} \in \mathbb{R}^{B \times S \times G \times C_t}.
\end{equation}

Operators are:
\begin{align}
\text{AND}(x) &= \prod_{i=1}^{C_t} x_i, \\
\text{OR}(x) &= 1 - \prod_{i=1}^{C_t}(1 - x_i), \\
\text{NOT}(x) &= 1 - x.
\end{align}

Applying to routed groups:
\begin{align}
a_{b,s,g} &:= \text{AND}(\operatorname{win}(u_t^{\text{AND}})_{b,s,g,:}), \\
o_{b,s,g} &:= \text{OR}(\operatorname{win}(u_t^{\text{OR}})_{b,s,g,:}), \\
n_{b,s,g,:} &:= \text{NOT}(\operatorname{win}(u_t^{\text{NOT}})_{b,s,g,:}).
\end{align}

AND and OR reduce to $[B,S,G]$.

NOT is applied elementwise and flattened to $[B,S,G \cdot C_t]$, to preserve
negation operation structure.

The concatenated representation has dimension:
\begin{equation}
G + G + G \cdot C_t = G(C_t + 2).
\end{equation}

\subsection{Layer Update and Residual Flow}

The concatenated operator outputs are projected:
\begin{equation}
\Delta L_t := \text{Linear}(\text{concat}(a, o, n)).
\end{equation}

The logic state is updated via:
\begin{equation}
L_t = \mathrm{RMSNorm}(L_{t-1} + \Delta L_t).
\end{equation}

\subsection{Output Fusion}

After the final layer, fusion is applied as in Eq.~\ref{eq:fusion}.
$\alpha$ may be fixed or learned, it directly controls
the contribution of the logic expert.

\section{Experiments}

\begin{table*}[tbp]
\centering
\normalsize
{\renewcommand{\arraystretch}{1.2}
\setlength{\tabcolsep}{6pt}
\begin{tabular}{|l|c|c|c|c|c|}
\hline
Model & Validation Accuracy & Validation Loss & Fusion $\alpha$ & Routing Entropy & Avg Grad Norm \\
\hline\hline
Baseline & 0.417 & 1.063 & -- & -- & 29.3765 \\
No-Gate Control & 0.525 & 0.820 & 0.00001 (fixed) & -- & 11.1387\\
Routed (intra, G=8) & 0.542 & 0.818 &  0.00001 (fixed) & 3.0675 & 10.5972 \\
Routed (inter, G=8) & 0.546 & 0.819 & 0.1 (fixed) & 3.0668 & 10.4154 \\
Routed (inter, G=8) & \textbf{0.562} & \textbf{0.816} & 0.01 (fixed) & 3.0666 & 10.4054 \\
Routed (inter, G=8) & 0.546 & 0.819 & 0.006 (learned) & 3.0666 & 10.5527 \\
Routed (inter, G=16) & 0.501 & 0.837 & 0.00001 (fixed) & 3.7483 & 10.9610 \\
\hline
\end{tabular}
}
\caption{Validation results. Abbreviations:
Baseline = backbone with task classifier head (no parallel branch); No-Gate Control = parameter-matched parallel branch without cross-attention/operator routing/compositional gating;
Routed = logic expert variants.}
\label{tab:main_results}
\end{table*}

\subsection{Setup}
We use Llama-3.1-8B \cite{grattafiori2024llama3herdmodels}
as the shared backbone for all variants.
We compare: (i) a Baseline
model, (ii) a No-Gate Control (parameter-matched), and
(iii+) the logic expert variants. The No-Gate Control serves as a
parameter-matched control to separate routing effects from pure
parameter-count gains.
All variants use the 3 label classifier head.
The original Llama head is not used.
In this control, the parallel branch excludes cross-attention, operator
routing, and compositional gating. All experiments use the AdamW
\cite{loshchilov2019decoupledweightdecayregularization}
optimizer and cross-entropy loss.
Values are collected on a single NVIDIA RTX PRO 6000 Blackwell Server
Edition GPU (CUDA 12.8) running on Google Colab.

\subsection{Logical Reasoning -- ProofWriter}
We use the ProofWriter \cite{tafjord2021proofwritergeneratingimplicationsproofs}
dataset to compare logical reasoning to the baselines.
ProofWriter is a logical reasoning dataset where each sample
contains a small rulebase (facts plus rules) and a query to evaluate as true,
false, or unknown. It is useful here because it directly tests multi-step
compositional inference, which matches the
intended behavior of the routed logic pathway.
We use deterministic, matched sampling across all models
(10k train, 5k validation, all problem depths).
Unless otherwise noted, all runs use matched seeds.
Main runs are trained for 3 epochs, longer
training only decreased classification due to overfitting.

\subsubsection{Beyond Logic -- Multi-Genre Natural Language (MNLI)}
To test performance beyond rulebase logic tasks, we evaluate
performance on Multi-Genre Natural Language (MNLI) \cite{N18-1101}.
MNLI allow us to evaluate the natural language understanding
of the model. MNLI has logical multi-class labels
(Entailment, Neutral, Contradiction),
this allows us to map to the same 3-label classification
head used for ProofWriter.

\subsection{Hyperparameter Evaluations}
Primary hyperparameters are logic dimension: $L$, gate count: $G$, gate mode:
(soft vs STE), gate window size: $C_t$,
window axis: (inter-token vs intra-token), fusion $\alpha$ mode:
(fixed vs learned), and initial fusion $\alpha$.
We validate on task validation sets, we use the No-Gate Control
architecture to avoid attributing gains to unmatched model size.

\subsection{Metrics}
Training metrics are validation accuracy and validation loss.
Operator mechanism metrics are routing entropy and fusion $\alpha$.
We also report average gradient norm (Avg Grad Norm) as an optimization-health
signal. Stable gradient norms generally indicate smoother updates,
very large norms can indicate unstable training dynamics, small
gradient norms indicates arrested training.

\section{Results}

\subsection{ProofWriter}
Table~\ref{tab:main_results} reports final validation accuracy
under the evaluation protocol used for each run, validation loss, and
fusion parameter state (fixed vs learned) after 3 epochs.
The parameter-matched control model substantially outperforms the
Baseline, indicating that additional parallel capacity alone
already improves performance, while routing variants
provide the strongest overall scores in this setting.
The best observed ProofWriter result is the inter-token routed variant with
$G=8$ and fixed $\alpha=0.01$ (0.562 val accuracy), compared with 0.525 for
No-Gate Control and 0.417 for Baseline. Inter-token routing is consistently
competitive with or better than intra-token routing at matched gate count,
while increasing to $G=16$ under these settings reduces accuracy (0.501),
indicating that more routed channels do not automatically improve quality.
Additional hyperparameters listed in the appendix.
The mechanism columns in Table~\ref{tab:main_results} show routed variants
have substantially lower gradient norms
(about 10.41--10.96) than the Baseline (29.38), while No-Gate Control is
intermediate (11.14). This is consistent with more stable update magnitudes in
the routed settings under the current training regime. Routing entropy is
tightly clustered around 3.067 for the $G=8$ routed runs, but rises to 3.7483
for $G=16$, where validation quality drops. In this sweep, higher routing
entropy at larger gate count appears to coincide with weaker task performance,
suggesting a more diffuse routing regime that does not translate into better
accuracy.

\subsection{Size and Transfer}

A key confound in models is added trainable parameters.
Table~\ref{tab:param_counts} confirms close parameter matching between No-Gate
Control and routed variants  (only small differences),
the single-task gains from the best routed model are
very unlikely to be explained solely by increased parameter-count.

\begin{table}[htbp]
\centering{
\begin{tabular}{|l|c|}
\hline
Model Variant & Trainable Parameters \\
\hline\hline
Baseline & 7,504,936,963 \\
No-Gate Control & 7,582,192,484 \\
Routed (G=8) & 7,579,068,164 \\
Routed (G=16) & 7,579,725,060 \\
\hline
\end{tabular}}
\caption{Trainable parameter counts for variants, inter/intra has no effect.}
\label{tab:param_counts}
\end{table}

\subsubsection{Multi-Genre Natural Language (MNLI)}
MNLI \cite{N18-1101}. combines sentence understanding with logic multi-class labels.
We map the labels to the ProofWriter classes (important for the next test):
contradiction $\rightarrow 0$, entailment $\rightarrow 1$, neutral $\rightarrow 2$.
For comparability, evaluation uses matched sampling. Notable hyperparameters
values in the appendix.

\begin{table}[htbp]
\centering
\resizebox{\linewidth}{!}{%
\begin{tabular}{|l|c|c|}
\hline
Model & MNLI Val Acc & MNLI Val Loss \\
\hline\hline
No-Gate Control & \textbf{0.89} & \textbf{0.34} \\
Routed (intra, G=8) &  0.88 & 0.36 \\
Routed (inter, G=8) & 0.88 & 0.39 \\
\hline
\end{tabular}
}
\caption{Performance on MNLI.}
\label{tab:mnli_transfer}
\end{table}

Table~\ref{tab:mnli_transfer} reports MNLI results. On MNLI,
No-Gate Control achieves the highest validation accuracy
(0.888), with routed variants close behind (0.880 intra, 0.878 inter). This
suggests broadly similar transfer performance at this scale; without
multi-seed variance estimates, these small margins should be interpreted as
comparable rather than decisive. Transfer behavior appears dependent on the
supervision regime and routing configuration.

\subsection{Multi-task Training}
For multi-task training, each mini-batch is sampled from a 50/50 mixture of
ProofWriter and MNLI. This experiment tests whether routing can jointly
optimize symbolic-style and open-domain NLI supervision without degrading
either task. Total training steps and optimizer hyperparameters are matched
to single-task runs, and we report per-task validation metrics after joint
training.

\begin{table}[htbp]
\centering
\resizebox{\linewidth}{!}{%
\begin{tabular}{|l|c|c|c|}
\hline
Model & ProofWriter Acc & MNLI Acc & Mean \\
\hline\hline
No-Gate Control & \textbf{0.544} & 0.879 & \textbf{0.712} \\
Routed (inter, G=8) & 0.536 & \textbf{0.881} & 0.709 \\
\hline
\end{tabular}
}
\caption{50/50 mixed ProofWriter and MNLI samples in the multi-task
training experiment.}
\label{tab:multitask_mix}
\end{table}

Table~\ref{tab:multitask_mix} reports the 50/50 multi-task evaluation.
In the joint setting, No-Gate Control attains the higher average accuracy
(0.712 vs 0.709), while Routed (inter, $G=8$) attains slightly higher MNLI
accuracy (0.881 vs 0.879). This indicates a small task trade-off under 50/50
mixing: routing is competitive in aggregate and can preserve transfer-side
performance, but these differences are very small and are best interpreted as
similar overall performance under the current multi-task configuration.

\section{Discussion}

\begin{figure*}[tbp]
\centering
\includegraphics[width=0.95\textwidth]{plots/routing_heatmap.png}
\caption{Routing-weight heatmap over transformer depth for a validation batch.
Rows are layers and columns are routed channels; color intensity indicates
mean routing weight.}
\label{fig:routing_heatmap}
\end{figure*}

The single-task results show that the routed logic expert variants can outperform both the
Baseline and the parameter-matched No-Gate Control, suggesting that gains are
not solely due to increased capacity and are influenced by compositional
routing bias. In contrast, the new 50/50 multi-task results show a tighter
regime with near-tied outcomes: No-Gate Control has a slightly higher average
accuracy (0.712 vs 0.709), while the routed variant is slightly higher on MNLI
within the joint setting (0.881 vs 0.879). Given the tiny margins, these
results are best read as similar performance pending multi-seed validation.
Single-task MNLI transfer follows the same pattern, while ProofWriter shows
larger routed gains. Overall, the evidence suggests a task-dependent benefit
profile in which routing is more helpful for logical-composition supervision.

Fusion $\alpha$ provides additional insight: models
with moderate nonzero $\alpha$ values indicate that the routing branch is
actively contributing to predictions, while near-zero values suggest the
model ignores the logic pathway. This highlights the importance of monitoring
mechanism-level metrics alongside task performance.

The inter-token capacity formulation changes the inductive bias of gate
composition from cross-channel local comparison to cross-token local
comparison. This makes each gate channel temporally local while preserving
channel independence. Empirically, this design should be interpreted together
with routing entropy and fusion $\alpha$: low entropy with stable nonzero $\alpha$
indicates specialized gates that remain active in final prediction, but these
mechanism diagnostics should be interpreted alongside task-level metrics and
held-out evaluation \cite{keysers2020measuringcompositionalgeneralizationcomprehensive}.

\begin{table}[t]
\centering
\small
\setlength{\tabcolsep}{4pt}
\begin{tabular}{|l|c|c|}
\hline
Model & Validation Acc & Validation Loss \\
\hline\hline
Routed-Gate-Only & \textbf{0.5348} & \textbf{0.832} \\
No-Gate Control & 0.4646 & 1.065 \\
\hline
\end{tabular}
\caption{ProofWriter evaluation for the earlier version without cross-attention, mean of 3 seeds evaluated.}
\label{tab:routed_gate_only_vs_nogate}
\end{table}

Table~\ref{tab:routed_gate_only_vs_nogate} suggests that cross-attention is
likely not the sole source of gains in the logic expert. However, these values
come from an earlier implementation with additional architectural and
hyperparameter differences, so both the baseline and logic-expert scores are
not directly comparable to the main results in Table~\ref{tab:main_results}.

Initialization of the learnable fusion $\alpha$ also affects optimization
stability. In our sweeps, values in the 0.01--0.25 range were generally stable,
with 0.1 yielding the strongest observed validation accuracy.

\subsection{Routing Visualization}

Figure~\ref{fig:routing_heatmap} visualizes routing specialization directly.
We see there the heatmap is sparse, with only a few hot-spots, this implies
there are several operators that dominate contribution.
This figure is used alongside routing entropy, which summarizes concentration
numerically, while the heatmap shows where specialization appears across layers.

\section{Limitations}

Additional tasks beyond the initial benchmark sets are
needed to establish any generalization.
We only evaluate one logical reasoning benchmark, more would allow us to
evaluate generalization, specifically of the logic expert.
The routing design also introduces additional hyperparameters
(gate count, window size, routing sparsity) whose
interactions are not yet fully explored.

\section{Conclusion}

This architecture shifts conditional inference from routing among
independent neural experts to routing over operator groups, yielding a
novel, compact, and interpretable compositional mechanism.
Across matched-capacity comparisons, results indicate that the routed logic
path can improve over both the backbone Baseline and a parameter-matched
No-Gate Control in single-task runs, with inter-token composition emerging as
a strong setting in our current sweeps. In MNLI transfer and joint 50/50
multi-task training, routed and No-Gate variants are near-parity under the
current setup.
Mechanism-level diagnostics (routing entropy and fusion
$\alpha$) provide direct visibility into whether routing is specialized and
whether the branch is actually used at prediction time.
This work establishes an implementation-complete foundation for operator-routed
reasoning. Immediate next steps are broader benchmark coverage, and targeted ablations.

{\small
\bibliographystyle{ieee}
\bibliography{egpaper_refs}
}

\appendix
\section*{Appendix}

\setcounter{figure}{0}
\setcounter{table}{0}
\setcounter{equation}{0}
\renewcommand{\thefigure}{A\arabic{figure}}
\renewcommand{\thetable}{A\arabic{table}}
\renewcommand{\theequation}{A\arabic{equation}}

\section{ProofWriter Evaluation Details}
This appendix summarizes the main ProofWriter configuration used in the
experiments.

\subsection{Settings and Hyperparameters}
\begin{itemize}
\item Batch size: 30 (train), 30 (eval)
\item Epochs: 3
\item Learning rate: $1\times10^{-5}$
\item Weight decay: 0.01
\end{itemize}

\subsection{Dataset(s)}
\begin{itemize}
\item ProofWriter Depth: \texttt{all}: (all allowed)
\item Train samples: 10,000
\item Validation samples: 5,000
\end{itemize}

\section{MNLI Evaluation Details}
For MNLI transfer, labels are remapped into the shared
three-class order used throughout this paper:
contradiction $\rightarrow 0$, entailment $\rightarrow 1$, and
neutral $\rightarrow 2$.

\subsection{Dataset(s)}
\begin{itemize}
\item Train samples: 4,000
\item Validation samples: 200
\end{itemize}

\subsection{Settings and Hyperparameters}
\begin{itemize}
\item Train batch size: 30
\item Eval batch size: 100
\item Learning rate: $1\times10^{-5}$
\item Weight decay: 0.01
\item Logic dimension / gates / heads: 256 / 8 / 4
\item Gate mode: \texttt{soft}
\end{itemize}


\section{Multi-task Evaluation Details}
For multi-task experiments, each batch is sampled from a mix
of ProofWriter and MNLI datasets.

\subsection{Dataset(s)}
\begin{itemize}
\item MNLI sampling ratio: 0.5
\item ProofWriter sampling ratio: 0.5
\item ProofWriter depth: \texttt{all}
\item Train/validation samples: 10,000 / 1,000
\end{itemize}

\subsection{Settings and Hyperparameters}
\begin{itemize}
\item Train batch size: 20
\item Eval batch size: 100
\item Epochs: 6
\item Learning rate: $1\times10^{-5}$
\item Weight decay: 0.01
\item Logic dimension / gates / heads: 256 / 8 / 4
\item Fusion $\alpha$ init: 0.1
\item Learn fusion $\alpha$: \texttt{True}
\item Gate axis: \texttt{inter\_token}
\item Gate mode: \texttt{soft}
\end{itemize}

\section{Routed-Gate-Only vs No-Gate Evaluation Details}
This appendix summarizes the overlapping configuration fields for
Table~\ref{tab:routed_gate_only_vs_nogate}.

\subsection{Settings and Hyperparameters}
\begin{itemize}
\item Epochs: 6
\item Learning rate: $1\times10^{-5}$
\end{itemize}

\section{ProofWriter Question}
This appendix section provides an example question from ProofWriter.

\subsection{Rulebase}
\begin{verbatim}
triple1: The cow is big.
triple2: The cow needs the dog.
triple3: The dog sees the rabbit.
triple4: The rabbit chases the cow.
triple5: The rabbit chases the dog.
triple6: The rabbit is big.
triple7: The rabbit sees the dog.
rule1: If the cow is blue and the cow needs the rabbit then the cow needs the dog.
rule2: If the cow chases the dog then the cow sees the rabbit.
rule3: If something is big then it chases the dog.
\end{verbatim}

\subsection{Question, Answer}
\begin{verbatim}
Question: The cow sees the rabbit?
Answer: true
\end{verbatim}

\end{document}
